\documentclass[UTF8,10pt,a4paper,twocolumn]{ctexart}
\usepackage[a4paper,left=15.5mm,right=15.5mm,top=16mm,bottom=17mm]{geometry}
\usepackage{xcolor,graphicx,booktabs,tabularx,array,caption}
\usepackage{amsmath,amssymb}
\usepackage{enumitem}
\usepackage{fancyhdr}
\usepackage{tikz}
\usepackage{hyperref}
\usetikzlibrary{positioning,arrows.meta,calc,fit}

\definecolor{lratblue}{RGB}{55,92,160}
\definecolor{papergray}{RGB}{243,245,247}
\definecolor{linegray}{RGB}{176,181,188}
\definecolor{accentred}{RGB}{190,61,55}
\definecolor{teal}{RGB}{43,124,111}
\definecolor{darkgray}{RGB}{55,58,63}

\hypersetup{colorlinks=true,linkcolor=lratblue,urlcolor=lratblue,citecolor=lratblue}
\graphicspath{{assets/}}
\setlength{\columnsep}{5.5mm}
\setlength{\parindent}{1em}
\setlength{\parskip}{1.5pt}
\setlength{\tabcolsep}{3.5pt}
\renewcommand{\arraystretch}{1.08}
\setlist{nosep,leftmargin=*}
\sloppy

\ctexset{
  section={format=\large\bfseries\color{lratblue},beforeskip=9pt,afterskip=3pt},
  subsection={format=\normalsize\bfseries\color{lratblue},beforeskip=7pt,afterskip=2pt},
  subsubsection={format=\small\bfseries\color{lratblue},beforeskip=5pt,afterskip=1pt}
}

\makeatletter
\fancypagestyle{mainpage}{%
  \fancyhf{}
  \fancyhead[C]{\small\textit{Agent 轨迹检索器的训练与验证}}
  \fancyhead[R]{\small\thepage}
  \fancyfoot{}
  \renewcommand{\headrulewidth}{0.35pt}
  \renewcommand{\footrulewidth}{0pt}
}
\fancypagestyle{firstpage}{%
  \fancyhf{}
  \fancyfoot[C]{\scriptsize\textit{内部研究稿 · v2.5 LRAT-style}}
  \renewcommand{\headrulewidth}{0pt}
  \renewcommand{\footrulewidth}{0pt}
}
\makeatother

\tikzset{
  pipe/.style={draw=lratblue!65,fill=lratblue!7,rounded corners=2pt,
    minimum width=2.35cm,minimum height=0.65cm,align=center,font=\footnotesize},
  pipegreen/.style={draw=teal!85!black,fill=teal!9,rounded corners=2pt,
    minimum width=2.35cm,minimum height=0.65cm,align=center,font=\footnotesize},
  pipeorange/.style={draw=accentred!80!black,fill=accentred!8,rounded corners=2pt,
    minimum width=2.35cm,minimum height=0.65cm,align=center,font=\footnotesize},
  pipearrow/.style={-{Latex[length=1.8mm]},draw=linegray,line width=0.65pt},
  metric/.style={draw=linegray,fill=papergray,rounded corners=2pt,
    minimum width=2.25cm,minimum height=0.58cm,align=center,font=\small}
}

\newcommand{\paperheader}{%
  \noindent
  \begin{minipage}[c]{0.60\textwidth}
    \textcolor{lratblue}{\sffamily\bfseries\large RUC}\hspace{1.5mm}
    \textcolor{darkgray}{\sffamily\small Renmin University of China}
  \end{minipage}%
  \hfill
  \begin{minipage}[c]{0.30\textwidth}
    \raggedleft\textcolor{lratblue}{\sffamily\bfseries CCIR}\\[-1mm]
    \raggedleft\scriptsize\textcolor{darkgray}{Retriever Study}
  \end{minipage}
  \par\vspace{1.5mm}\hrule height 0.35pt\vspace{3mm}
}

\newcommand{\abstractbox}[1]{%
  \noindent
  \colorbox{papergray}{%
    \parbox{0.94\textwidth}{\small #1}%
  }\par
}

\newcommand{\tightcaption}[1]{%
  \refstepcounter{figure}%
  \par\vspace{1mm}\noindent{\small\textbf{图 \thefigure \,|\,} #1}\par\vspace{1mm}
}

\begin{document}
\onecolumn
\thispagestyle{firstpage}
\paperheader
\begin{center}
  {\LARGE\bfseries Agent 轨迹检索器的训练与验证}\\[1mm]
  {\large\textit{一条可追溯、query 隔离、可独立验收的训练链路}}\\[2mm]
  \textbf{DefaultGroup\textsuperscript{1}}\\[-0.5mm]
  \small 中国人民大学高瓴人工智能学院，CCIR 参赛队
\end{center}
\vspace{2mm}
\abstractbox{\textbf{摘要：}
传统检索器主要从人的查询、点击和停留日志中学习，而 Agent 会在多轮 Think、Search、Browse 和 Reason 中提出中间 query，并把检索结果作为下一轮推理的观察输入。LRAT 论文说明，Search$\rightarrow$Browse 行为、未 Browse 候选和 post-browse reasoning 可以提供与文档效用相关的训练信号。本文围绕官方已发布 training pairs、Agent trajectories、离线 corpus 和原始 Qwen3-Embedding-0.6B，建立一条不重新生成 trajectory、但可追溯、可隔离、可训练、可验收的工程链路。我们回溯 pair 字段并完成 stable/ambiguous/mismatch 审计，按 normalized query 整组切分 train/dev/locked test，采用 1 positive + 5 negatives 的加权对比训练，并用开发集门控学习率与独立冻结验收。最终 A 榜为 Total 40.6、Recall 46.3、Success 21.1、Avg\_Steps 16.0；三种 Agent 的现有 B 榜评测记录平均 Total 为 56.18。}
\vspace{2mm}
\begin{center}
\begin{tikzpicture}[node distance=0.28cm and 0.22cm]
  \node[pipe] (a) {固定官方输入};
  \node[pipegreen,right=of a] (b) {轨迹 / pair\\来源审计};
  \node[pipe,right=of b] (c) {normalized query\\整组隔离};
  \node[pipegreen,right=of c] (d) {1+5 加权\\对比训练};
  \node[pipeorange,right=of d] (e) {A/B 评测\\边界解释};
  \node[pipegreen,right=of e] (f) {冻结、manifest\\独立回读};
  \foreach \x/\y in {a/b,b/c,c/d,d/e,e/f}{\draw[pipearrow] (\x) -- (\y);}
\end{tikzpicture}
\par\vspace{1mm}
\begin{tikzpicture}[node distance=0.24cm and 0.25cm]
  \node[metric] (m1) {\textbf{40.6}\\[-1mm]\scriptsize A Total};
  \node[metric,right=of m1] (m2) {\textbf{46.3}\\[-1mm]\scriptsize Recall};
  \node[metric,right=of m2] (m3) {\textbf{21.1}\\[-1mm]\scriptsize Success};
  \node[metric,right=of m3] (m4) {\textbf{56.18}\\[-1mm]\scriptsize B 平均 Total};
\end{tikzpicture}
\tightcaption{本文的训练与证据链路。改进集中在监督来源、query 隔离、训练控制和独立交付，而不是重新提出 LRAT 的轨迹生成机制。}
\end{center}

\clearpage
\pagestyle{mainpage}
\twocolumn

\section{引言}

信息检索系统长期围绕人的信息需求设计：用户提交一个相对稳定的 query，系统返回结果页，点击和停留等交互再被用作弱监督。这个假设在 Agentic Search 中发生了变化。深度研究 Agent 会先根据当前 reasoning 产生中间 query，观察 Search 返回的候选摘要，再决定是否 Browse 某篇全文；全文证据又会改变下一轮 reasoning、query 和最终答案。检索结果不再是服务终点，而是 Agent 的感知输入。

这种使用方式带来一个具体的训练问题：一个文档是否“相关”，不能只由它与最终问题的主题相似度决定，还要看它能否推进 Agent 当前的局部推理。若把面向人的通用相关性或未经审计的扁平 pair 直接用于训练，就难以回答三个问题：正负样本为什么成立，同一 query 是否泄漏到验证集，以及一个训练结果能否由独立过程重新加载和回读。

LRAT 将这一问题表述为从 Agent trajectory 学习检索器\cite{lrat}。它从 Search$\rightarrow$Browse 转移中获取正例候选，把同一候选集内未 Browse 的文档作为经验负例，并用 Browse 后 reasoning 的长度表达相关性强度。本文不把 LRAT 的轨迹生成、LLM judge 或论文实验重新包装成本项目的新贡献；我们的目标是围绕固定官方输入，把这套监督变成一条可审计、可控且可复核的训练链路。

\textbf{本文的实际贡献如下：}
\begin{enumerate}
  \item 将官方 pair 的 query、正负文档和 \texttt{reweight\_rate} 逐条回溯到 trajectory 事件和 corpus，并显式区分 stable、ambiguous 和 mismatch。
  \item 对 query 做 strip、连续空白合并和 Unicode casefold，按完整 normalized-query group 隔离 train、dev 和 locked test。
  \item 固定 1 positive + 5 negatives、128/512 token 预算、官方权重和 cross-device negatives，记录真实训练配置与 loss 归一化方式。
  \item 用 dev1500 进行配置门控，随后独立加载、冻结、生成 manifest，并把 A 榜、B 榜和交付状态分开解释。
\end{enumerate}

\section{相关工作}

\textbf{面向人的检索训练。}
传统 dense retriever 和 learning-to-rank 方法通常以人工相关性、点击或其他用户交互作为监督。它们适合“query $\rightarrow$ 结果页”的服务形式，但这些信号并不直接描述 Agent 在多轮推理中为什么选择某篇全文，也不记录一次检索是否改变了后续行动。

\textbf{从 Agent 轨迹学习。}
LRAT 的核心观察是，Agent 的 Search、Browse 和 post-browse reasoning 构成了一种不同于人类点击日志的行为监督。Browse 表示 Agent 愿意读取全文，未 Browse 候选在论文设定中提供了经验负例，而 reasoning length 可以提供强度信息。本文沿用这套问题语言，用它解释官方 pair 的字段；同时严格区分“论文提出的训练范式”和“本项目在固定官方 pair 上完成的工程实现”。

\textbf{本文的定位。}
本项目不是对 LRAT 论文的完整再现，也没有重新运行论文中的 10K query、30B Agent、LLM judge 或 data flywheel。我们研究的是一个受竞赛合同约束的 retriever 构建过程：原始模型、官方 pairs 和推理 contract 是固定输入，改进点是 provenance、query 隔离、训练控制和交付验收。

\section{预备知识}

\subsection{Agent trajectory}

设初始任务为 $q$。在第 $t$ 轮，Agent 维护 reasoning state $r_t$，先产生一个中间 Search query $q_t$，检索系统返回候选文档集合 $\mathcal{D}_t=\{d_{t,i}\}_{i=1}^{K}$ 及摘要。Agent 随后可以选择 Browse 某个候选文档，获得全文证据并把它写回新的 reasoning state。一个简化的交互过程为
\[
\text{Think}\rightarrow\text{Search}\rightarrow\text{Browse}\rightarrow\text{Reason}\rightarrow\text{Answer}.
\]
因此，训练信号需要描述“候选文档被如何使用”，而不仅是 query 和文档的静态相似度。

\subsection{任务定义与证据边界}

本项目的输入是官方已发布的 training pairs、原始 Qwen3-Embedding-0.6B、轨迹与离线 corpus；训练目标是得到一个可独立加载的 dense retriever，并在官方 A/B evaluator contract 下读取 query/document 表示。我们不把 dev1500 当作正式泛化证明，也不使用 locked test 进行学习率或 checkpoint 选择。A 榜是固定 Agent 的端到端证据，B 榜只说明在现有 Agent 评测记录上的可运行性，不能外推为所有 Agent 的普遍保证。

\section{数据审计与 query 隔离}

\subsection{Pair provenance}

官方 pair 行包含 query、正例文档引用、负例文档列表和 \texttt{reweight\_rate}。我们用 \texttt{build\_provenance.py} 将这些字段回到 Search/Browse/reasoning 事件，再用 \texttt{preprocess.py} 检查字段类型、负例 ID 和 corpus 文本一致性。全量 96,504 条 pair 中，96,366 条能够稳定回溯，138 条为 ambiguous，0 条 mismatch；1,989,015 个文档引用和 959,042 个唯一文档 ID 均无缺失或文本不一致。

这里需要保留一个版本边界：LRAT PDF 的一处方法描述写出 91,713 条 pairs，而本地官方发布文件有 96,504 行。两者属于不同发布或处理版本，本文只使用当前项目锁定的 96,504 行，不把两个数字拼成一个训练集，也不据此声称完整复现论文数据规模。

\subsection{Query-disjoint split}

如果按 source row 随机拆分，同一 query 的多条 pair 可能同时落入 train 和 dev，dev 指标就会混合记忆重复 query 与真正的 query 适应能力。我们先做去首尾空格、合并连续空白和 Unicode casefold，再把同一 normalized query 的全部 rows 作为一个 group，按确定性顺序整组分配。

\begin{table*}[t]
\centering
\caption{Query-disjoint 数据合同。}
\label{tab:split}
\small
\begin{tabular}{lrrp{6.3cm}}
\toprule
集合 & Query groups & Source rows & 用途\\
\midrule
Train & 78,890 & 94,113 & 全参数训练\\
Dev & 1,500 & -- & 学习率与配置选择\\
Locked test & 500 & -- & 最终诊断，不回流选择\\
\bottomrule
\end{tabular}
\end{table*}

最终 train/dev/test 的 normalized-query overlap 为 0。Dev 只承担开发与门控职责；locked test 没有用于当前配置选择。这一分割提升的是证据边界的清晰度，而不是自动保证榜单得分提高。

\section{训练方法}

\subsection{Training sample construction}

每条官方 pair 行构造成一个 query、一个 positive 和从该行 \texttt{neg} 列表抽取的五个 negatives，组成六路候选组，positive 固定在组内位置 0。query 和 passage 上限分别为 128 和 512 tokens；它们是控制六路候选显存与计算量的训练预算，不被表述为理论无损压缩。榜单推理仍遵守官方 evaluator contract。

\subsection{Encoder and weighted contrastive learning}

模型从原始 Qwen3-Embedding-0.6B 开始全参数更新，使用共享 encoder、last-token pooling、L2 normalization 和 1024 维向量。对 query $q$ 和 passage $d$，相似度为
\begin{equation}
s(q,d)=\frac{q^{\mathsf T}d}{\tau},\qquad \tau=0.02.
\end{equation}
设第 $j$ 个 query 的对比交叉熵为 $\ell_j$，官方权重为 $w_j$，当前实现的 micro-batch 目标为
\begin{equation}
\mathcal{L}=\frac{\sum_j w_j\ell_j}{\sum_j w_j}.
\end{equation}
权重改变 query 对总 loss 的贡献，不改变正负标签。两张 A40 各处理 1 个 query；cross-device gather 后，另一 query 的候选文档成为 in-batch negatives。gradient accumulation=4 只累积梯度，不合并四个 micro-step 的候选池，effective query batch 为 8。

\subsection{Configuration gate}

在固定 dev1500 上比较四个 500-step 学习率点。$2\times10^{-6}$ 同时取得最高 R@1 和 MRR，才进入完整一轮训练。完整配置为 1 epoch、11,764 optimizer steps、per-device batch 1、accumulation 4、bf16、gradient checkpointing 和 cross-device negatives；不使用外部语料、教师模型、LoRA 或 Adapter。

\begin{table}[t]
\centering
\caption{Dev1500 上的短训学习率选择。}
\label{tab:lr}
\scriptsize
\begin{tabular}{lrr}
\toprule
Learning rate & R@1 & MRR\\
\midrule
$3\times10^{-7}$ & 0.4913 & 0.6419\\
$5\times10^{-7}$ & 0.5073 & 0.6571\\
$1\times10^{-6}$ & 0.5433 & 0.6831\\
\textbf{$2\times10^{-6}$} & \textbf{0.5640} & \textbf{0.6993}\\
\bottomrule
\end{tabular}
\end{table}

\section{实验}

\subsection{Experimental setup}

评测使用固定官方 evaluator。训练输入为 94,113 行 query-disjoint train，基础模型为原始 Qwen3-Embedding-0.6B，硬件为 2$\times$ NVIDIA A40。我们保留原始官方 \texttt{reweight\_rate}，不加入外部数据或教师模型。正式结果以 A/B 榜为准，dev 指标只负责配置选择。

\begin{table}[t]
\centering
\caption{最终训练配置。}
\label{tab:config}
\scriptsize
\begin{tabular}{ll}
\toprule
项目 & 配置\\
\midrule
Base model & Qwen3-Embedding-0.6B\\
Train rows & 94,113\\
GPU / epoch & 2$\times$A40 / 1\\
Optimizer steps & 11,764\\
Batch / accumulation & 1 / 4\\
Group / temperature & 6 / 0.02\\
Query / passage cap & 128 / 512\\
Precision / negatives & bf16 / cross-device\\
\bottomrule
\end{tabular}
\end{table}

\subsection{A leaderboard}

\begin{table}[t]
\centering
\caption{固定 Agent 的 A 榜端到端结果。}
\label{tab:aleaderboard}
\scriptsize
\begin{tabularx}{\columnwidth}{@{}>{\raggedright\arraybackslash}Xrrrr@{}}
\toprule
候选 & Total $\uparrow$ & Recall $\uparrow$ & Success $\uparrow$ & Avg\_Steps $\downarrow$\\
\midrule
Qwen3-Embedding-0.6B (Base) & 28.5 & 28.5 & 10.6 & 17.8\\
官方 pairs 一轮 & 40.1 & 44.8 & 21.0 & 15.4\\
Query-disjoint ($1\times10^{-6}$) & 39.5 & 45.0 & 19.8 & 15.9\\
\textbf{Query-disjoint ($2\times10^{-6}$)} & \textbf{40.6} & \textbf{46.3} & \textbf{21.1} & 16.0\\
\bottomrule
\end{tabularx}
\end{table}

最终候选相对 Base 提升 12.1 个 Total 点；在同一 query-disjoint 配方中，相对 $1\times10^{-6}$ 提升 1.1 个 Total 点。相对官方 pairs 一轮，Recall 增加 1.5、Success 增加 0.1，但 Avg\_Steps 增加 0.6，说明收益并非所有指标单调改善。因而本文把结果解释为一个受控配置在固定 Agent 上的端到端验证，而不是单一 dev 指标的泛化承诺。

\subsection{B leaderboard and delivery}

现有 B 榜评测记录中，三种 Agent 的 Total 为 67.42、60.96 和 40.16，平均 56.18。该结果支持跨 Agent 可运行性，但分数差异也说明 Agent 自身的模型能力和行动策略会影响端到端表现，不能外推为普遍保证。

交付路径遵循“checkpoint $\rightarrow$ freeze + manifest $\rightarrow$ 独立 CPU load/validator $\rightarrow$ HF、GitHub、ZIP”的顺序。公开模型、代码仓库和 reproduction package 的对象身份可以回读；这些准备与验收不等于自动完成比赛平台上传，平台提交仍是独立的授权门。

\subsection{Limitations and evidence boundary}

第一，dev1500 不是 held-out leaderboard，也不能证明对未知 Agent 的泛化；locked test 未用于当前选择。第二，M10 采用官方已发布 pairs 和当前项目训练配置，不能写成 LRAT 论文完整实验复现。第三，训练阶段 128/512 的 truncation contract 与榜单推理 contract 分开；它控制训练成本，不等于评测输入被同样截断。第四，A/B 榜回答不同问题：A 榜固定 Agent 的综合效果，B 榜观察多个 Agent 的现有评测表现。

\section{结论}

本文按照 LRAT 的问题组织方式，从“Agent 如何使用检索结果”出发，经过相关工作、轨迹预备知识、监督审计、训练方法和实验边界，最后回到可复核结论。我们的实际改进不是重新提出一个新的 trajectory mining 范式，而是把固定官方监督转化为一条可相信的工程链路：来源可追溯，query 可隔离，训练可控制，模型可独立加载，结果可按 A/B 榜分别解释。

更具体地说，来源追溯回答监督从哪里来，query-disjoint 回答验证是否泄漏，1+5 加权训练回答如何在有限预算内保留强度信息，dev gate 和 freeze/manifest 回读回答结果是否仍是同一个对象。当前证据支持 A 榜 Total 40.6 和 B 榜现有记录平均 56.18；它同时要求我们保留“不是完整 LRAT 复现、不是普遍 Agent 泛化、不是所有指标均改善”的边界。

\appendix
\section{Evidence and reproducibility}

本稿正文只保留能帮助理解逻辑的数字；完整身份和交付证据见项目附录 \texttt{05\_复现验收附录.md}。当前关键对象包括：最终 M10 权重 SHA 前缀 \texttt{cea87cca...a9984}，公开模型 commit \texttt{9cdfdba...16bb}，canonical GitHub \texttt{main@31949f6...7bba}，以及 B 榜三字段 JSONL SHA 前缀 \texttt{02c27117...b3f49}。这些摘要用于定位对象，不替代完整 manifest。

\begin{table}[t]
\centering
\caption{可回读证据层。}
\label{tab:evidence}
\scriptsize
\begin{tabular}{p{1.25cm}p{6.0cm}}
\toprule
层级 & 检查内容与当前边界\\
\midrule
输入 & manifest 锁定 Base model、raw pairs、provenance、corpus 与 split；不引入外部语料\\
加载 & 新 CPU 进程执行 query/document smoke；证明对象可加载，不证明不同硬件训练逐字节一致\\
训练 & 回读命令、配置、checkpoint 和 loss 归一化；证明当前实现可解释，不等于复现论文全部实验\\
交付 & 回读 HF、GitHub、ZIP 的 manifest 与 validator；证明可复核交付，不等于比赛平台已上传\\
\bottomrule
\end{tabular}
\end{table}

\begin{thebibliography}{9}
\bibitem{lrat}
Yuqi Zhou, Sunhao Dai, Changle Qu, Liang Pang, Jun Xu, and Ji-Rong Wen.
\newblock Learning to Retrieve from Agent Trajectories.
\newblock Local LRAT paper artifact; official repository:
\url{https://github.com/Yuqi-Zhou/LRAT}.

\bibitem{pairs}
Yuqi-Zhou.
\newblock LRAT training pairs and official project data contract.
\newblock Hugging Face dataset:
\url{https://huggingface.co/datasets/Yuqi-Zhou/LRAT-Train}.

\bibitem{qwen}
Qwen Team.
\newblock Qwen3-Embedding-0.6B model card.
\newblock Hugging Face:
\url{https://huggingface.co/Qwen/Qwen3-Embedding-0.6B}.
\end{thebibliography}

\end{document}
